%! Subsubsección: Grafos: la sociedad como red
% Cuerpo sin encabezado de sección.

Una red social puede representarse formalmente como un grafo $G = (V, E)$, donde el conjunto de nodos $V$ corresponde a los individuos que participan en el sistema y el conjunto de aristas $E \subseteq V \times V$ representa los vínculos entre ellos. Según la naturaleza de dichos vínculos, el grafo puede ser \emph{no dirigido} (relaciones simétricas de amistad o conocimiento mutuo), \emph{dirigido} (relaciones asimétricas como el seguimiento en redes digitales), \emph{ponderado} (cuando se asigna un valor numérico de confianza o intensidad a cada arista) o \emph{multicapa} (cuando coexisten distintos tipos de red sobre el mismo conjunto de individuos). Esta representación unificada permite aplicar al estudio de la influencia social el aparato matemático de la teoría de grafos y la ciencia de redes \cite{AlbertBarabasi2002}.

\paragraph{Métricas estructurales.}
La caracterización cuantitativa de una red requiere el cálculo de un conjunto de métricas estructurales que el módulo \texttt{compute\_graph\_metrics} del proyecto evalúa sobre cada topología generada. El \emph{grado} $k_i$ de un nodo $i$ es el número de aristas que lo inciden; en redes dirigidas se distinguen el grado de entrada $k_i^{\text{in}}$ y el de salida $k_i^{\text{out}}$. La distribución de grado $P(k)$ describe la fracción de nodos con grado $k$ y constituye la huella estructural más característica de cada modelo de red. El \emph{coeficiente de clustering} $C_i$ mide la densidad de triángulos en el vecindario de $i$, mientras que la \emph{longitud media de camino} $\langle \ell \rangle$ refleja la eficiencia global de la propagación de información a través de la red. Las \emph{centralidades} de intermediación, de vector propio y PageRank cuantifican, respectivamente, el control de flujo, la influencia transitiva y la notoriedad de cada nodo. La \emph{asortatividad} indica si los nodos de grado elevado tienden a conectarse entre sí o con nodos periféricos.

Cuando la distribución de grado sigue una \emph{ley de potencias}, se cumple

\begin{equation}
  P(k) \propto k^{-\gamma},
  \label{eq:powerlaw}
\end{equation}

donde el exponente $\gamma$ se estima mediante el estimador de máxima verosimilitud de Clauset, Shalizi y Newman \cite{Clauset2009}:

\[
  \hat{\gamma} = 1 + n \left[ \sum_{i=1}^{n} \ln\!\left( \frac{k_i}{k_{\min} - \tfrac{1}{2}} \right) \right]^{-1},
\]

con $k_{\min}$ el grado mínimo a partir del cual se ajusta la ley y $n$ el número de nodos con $k_i \geq k_{\min}$. Este estimador es asintóticamente insesgado y se utiliza en el código para caracterizar automáticamente las redes de escala libre generadas.

\paragraph{Topologías implementadas.}
El simulador incorpora tres modelos de red canónicos que sirven de base a todas las configuraciones de experimento.

El modelo de \emph{Erdős–Rényi} $\mathcal{G}(n, p)$ \cite{ErdosRenyi1960} construye el grafo conectando cada par de nodos de forma independiente con probabilidad $p$. El grado de cada nodo sigue una distribución de Poisson con media $\langle k \rangle = (n-1)p$. Esta topología se usa exclusivamente como \emph{baseline} aleatorio: carece de hubs, presenta clustering bajo y no reproduce las heterogeneidades observadas en redes humanas reales.

El modelo \emph{libre de escala} \cite{BarabasiAlbert1999, AlbertBarabasi2002} se basa en el mecanismo de adjunción preferencial: los nodos nuevos se enlazan con mayor probabilidad a nodos ya bien conectados, lo que genera la distribución de grado~\eqref{eq:powerlaw} con exponente $\gamma \in [2, 3]$ y una minoría de nodos con grado excepcional (\emph{hubs}). En el código se utiliza la generalización dirigida de Bollobás et al.\ \cite{Bollobas2003}, implementada mediante \texttt{nx.scale\_free\_graph}, que introduce tres parámetros de mezcla $\alpha + \beta + \gamma = 1$: $\alpha$ controla la probabilidad de añadir un nodo nuevo con arista de entrada, $\beta$ la de añadir una arista entre nodos existentes, y $\gamma$ la de añadir un nodo nuevo con arista de salida; los sesgos $\delta_{\text{in}}$ y $\delta_{\text{out}}$ modulan la preferencia de adjunción. Esta topología modela la \emph{red digital}: la distancia media entre usuarios en Facebook se aproxima a 3{,}5 saltos \cite{Backstrom2012}, confirmando el efecto del mundo pequeño inducido por los hubs.

El modelo de \emph{Watts–Strogatz} \cite{WattsStrogatz1998} parte de un anillo regular donde cada nodo está conectado a sus $k$ vecinos más próximos y, posteriormente, recablea cada arista con probabilidad $p$. Para valores intermedios de $p$ se obtiene simultáneamente un coeficiente de clustering elevado (herencia de la estructura regular) y caminos medios cortos proporcionales a $\log n$ (consecuencia del recableado aleatorio). Esta topología modela la \emph{red analógica}, cuya estructura refleja la organización comunitaria de las relaciones presenciales.

La figura~\ref{fig:pk} compara las distribuciones de grado de los tres modelos en escala logarítmica.

\begin{figure}[h!]
  \centering
  \fbox{\parbox{0.85\linewidth}{\centering\vspace{1.5cm}
    PENDIENTE (figura a generar): Distribuciones de grado $P(k)$ en escala log-log para los tres modelos de red implementados: Erdős–Rényi (distribución de Poisson, caída exponencial), libre de escala (recta de pendiente $-\gamma$, hubs visibles) y Watts–Strogatz (distribución acotada con pico estrecho). Ejes: $\log k$ frente a $\log P(k)$.
    \vspace{1.5cm}}}
  \caption{Distribuciones de grado $P(k)$ en escala log-log para los modelos Erdős–Rényi, libre de escala (Barabási–Albert) y Watts–Strogatz. La recta de pendiente $-\gamma$ en la red libre de escala evidencia el comportamiento de ley de potencias~\eqref{eq:powerlaw}.}
  \label{fig:pk}
\end{figure}

\paragraph{Confianza por capa.}
La confianza entre dos agentes no es homogénea ni estática, sino que depende de la capa de red en la que se establece el vínculo. El módulo \texttt{src/graphs/trust.py} implementa dos modelos diferenciados.

En la \emph{capa analógica}, la confianza se asigna conforme al modelo de Dunbar \cite{Dunbar1992}: la distancia circular entre dos nodos $u$ y $v$ en el anillo de Watts–Strogatz se define como $d = \min(|u - v|,\, n - |u - v|)$, y el valor de confianza se asigna en función de la banda de Dunbar a la que pertenece dicha distancia — núcleo íntimo ($\leq 5$), círculo cercano ($\leq 15$), amigos ($\leq 50$) o conocidos ($\leq 150$) —, de modo que a mayor distancia corresponde menor confianza. Esta graduación reproduce la estructura jerárquica observada empíricamente en las redes de relación humana presencial.

En la \emph{capa digital}, la confianza se modela mediante una distribución log-normal cuya media y varianza dependen de las propiedades estructurales de emisor y receptor:

\begin{equation}
  \tau^{\text{dig}}_{a \to b} = \mathrm{clip}\!\left(\,\mathcal{N}(\mu, \sigma) + 0{,}04 \cdot \min(c_{ab}, 5),\; 0,\; 1\right),
  \label{eq:digital-trust}
\end{equation}

donde $\mu = 1 / (1 + 0{,}005 \cdot k_b^{\text{in}})$, $\sigma = 0{,}05 + 0{,}003 \cdot k_a^{\text{out}}$, $k_b^{\text{in}}$ es el in-grado del receptor $b$, $k_a^{\text{out}}$ el out-grado del emisor $a$, y $c_{ab}$ el número de vecinos comunes entre $a$ y $b$. La interpretación es directa: los hubs con muchos seguidores generan una confianza por contacto decreciente (efecto de \emph{saturación}), mientras que compartir vecinos introduce un pequeño bono de credibilidad basado en el principio de cierre triádico.

\paragraph{Redes multicapa.}
Cuando se desea modelar simultáneamente las interacciones analógicas y digitales de una misma población, el formalismo de redes multicapa \cite{Kivela2014, DeDomenico2013} proporciona el marco apropiado. Una red multicapa se define formalmente como la tupla $\mathcal{M} = (V_\mathcal{M}, E_\mathcal{M}, V, L)$, donde $L = \{L_1, \ldots, L_d\}$ es el conjunto de capas, $V_\mathcal{M} \subseteq V \times L$ el conjunto de nodos-capa y $E_\mathcal{M}$ incluye tanto aristas \emph{intra-capa} (dentro de una misma capa) como aristas \emph{inter-capa} (entre capas).

En el simulador, la clase \texttt{MultiLayerGraph} construye un \texttt{MultiDiGraph} de NetworkX con dos capas: la digital (libre de escala, prioritaria) y la analógica (Watts–Strogatz). Se mantiene el invariante de que ningún par $\{u, v\}$ aparece simultáneamente en ambas capas, evitando duplicidades que distorsionarían la dinámica de propagación. La \emph{participación} de un nodo en las distintas capas se cuantifica mediante el índice

\begin{equation}
  P_i = 1 - \sum_{\ell} \left( \frac{k_{i,\ell}}{k_i} \right)^2,
  \label{eq:participacion}
\end{equation}

donde $k_{i,\ell}$ es el grado del nodo $i$ en la capa $\ell$ y $k_i = \sum_\ell k_{i,\ell}$ su grado total. Un valor $P_i \approx 1$ indica participación equilibrada entre capas, mientras que $P_i \approx 0$ señala que toda la actividad del nodo se concentra en una única capa. Este índice, junto con la correlación de grado entre capas, permite diagnosticar la dualidad analógico-digital de cada agente.

\begin{figure}[h!]
  \centering
  \fbox{\parbox{0.85\linewidth}{\centering\vspace{1.5cm}
    PENDIENTE (figura a generar): Esquema de la red multicapa analógico-digital. Se muestran dos planos paralelos: el superior representa la capa digital (libre de escala, con hubs bien visibles) y el inferior la capa analógica (Watts–Strogatz, estructura de anillo con conexiones locales). Las aristas inter-capa (punteadas verticales) conectan el mismo nodo en ambas capas. El código de colores distingue aristas intra-capa analógica (azul), intra-capa digital (naranja) e inter-capa (gris).
    \vspace{1.5cm}}}
  \caption{Arquitectura de la red multicapa analógico-digital. La capa superior (digital, libre de escala) y la capa inferior (analógica, Watts–Strogatz) comparten el mismo conjunto de nodos $V$; las aristas verticales representan la correspondencia inter-capa.}
  \label{fig:multicapa}
\end{figure}

\paragraph{Redes reales (SNAP).}
Además de las topologías sintéticas, el simulador permite cargar redes empíricas del repositorio SNAP \cite{LeskovecKrevl2014}. Los conjuntos de datos utilizados — redes de ego de Facebook y Twitter, y la red de confianza de Epinions — exhiben los sesgos estructurales característicos de las redes humanas: distribución de grado de cola pesada, alto clustering local y efecto de mundo pequeño. La clase \texttt{SNAPGraph} se encarga de descargar y adaptar estos datos al formato interno del simulador; cuando el grafo original es no dirigido, la orientación de las aristas se asigna de forma probabilística ponderada por el grado de los nodos extremos, aproximando la asimetría típica de las relaciones de influencia en entornos digitales.

\paragraph{Comparativa analógica-digital.}
La tabla~\ref{tab:analog-digital} resume las diferencias estructurales más relevantes entre la red analógica y la digital en el contexto del simulador.

\begin{table}[h!]
  \caption{Comparativa estructural entre la red analógica (Watts–Strogatz) y la red digital (libre de escala) en el simulador.}
  \label{tab:analog-digital}
  \centering
  \begin{tabular}{@{}l p{5cm} p{5.8cm}@{}}
    \toprule
    \textbf{Dimensión} & \textbf{Red analógica} & \textbf{Red digital} \\
    \midrule
    Modelo de base & Watts–Strogatz & Barabási–Albert dirigido \cite{Bollobas2003} \\
    Tamaño de la vecindad & Limitado (bandas de Dunbar \cite{Dunbar1992}) & Ilimitado; hubs con $k \gg \langle k \rangle$ \\
    Coste de vínculo & Alto (tiempo, proximidad física) & Bajo (un clic, coste casi nulo) \\
    Distribución de grado & Acotada, pico estrecho & Ley de potencias, $\gamma \in [2,3]$~\eqref{eq:powerlaw} \\
    Distancia media $\langle \ell \rangle$ & $\mathcal{O}(\log n)$, relativamente larga & Muy corta ($\approx 3{,}5$ saltos \cite{Backstrom2012}) \\
    Coeficiente de clustering & Alto (estructura comunitaria) & Bajo (hubs con vecindarios dispersos) \\
    Mediación algorítmica & Ninguna & Presente (recomendaciones, \emph{trending}) \\
    Confianza por arista & Según bandas de Dunbar & Log-normal modulada por grado~\eqref{eq:digital-trust} \\
    \bottomrule
  \end{tabular}
\end{table}
